\section{The classical product: explicit remainder and a rational certificate}
\label{app:product}

The product in \eqref{eq:critical-product} is the established
Poisson exponential-functional transform of
\citet{BertoinBianeYor2004}. The calculations here give an explicit
truncation error and a useful large-argument form for this particular
normalization. They concern the conditional size argument $q$,
not the calendar-time tail $h(t)$.

\begin{proposition}[Periodic large-size correction]
\label{prop:product-periodic}
For $q\geq1$, write $\log_2q=n+\theta$ with integer $n\geq0$ and
$0\leq\theta<1$; the hypothesis $q\geq1$ is what makes $n\geq0$. Then
\begin{equation}
 \log\Psi(q)=-\frac{(\log q)^2}{2\log2}
                +\frac12\log q-\mathcal P(\theta)+\mathcal R(q),
 \label{eq:product-periodic}
\end{equation}
where
\begin{align}
 \mathcal P(\theta)
 &=\frac{\log2}{2}(\theta-\theta^2)
   +\sum_{j=0}^\infty\log(1+2^{-j-\theta})
   +\sum_{\ell=1}^\infty\log(1+2^{\theta-\ell}),
 \label{eq:product-periodic-function}\\
 \mathcal R(q)&=\sum_{r=0}^\infty\log(1+q^{-1}2^{-r}),
             \qquad 0\leq\mathcal R(q)\leq2/q.
 \label{eq:product-remainder}
\end{align}
The function $\mathcal P$ extends to a bounded Lipschitz
one-periodic function on $\R$, with Lipschitz constant at most
$\tfrac92\log2$. More precisely,
\begin{equation}
 \frac2q-\frac{2}{3q^2}\leq\mathcal R(q)
       \leq\frac2q-\frac{2}{3q^2}+\frac8{21q^3}.
 \label{eq:product-refined-remainder}
\end{equation}
The bounds \eqref{eq:product-remainder} and
\eqref{eq:product-refined-remainder} hold for every $q>0$.
\end{proposition}

\begin{proof}
Put $a=n+\theta=\log_2q$. Split $-\log\Psi(q)$ at $k=n$ and
factor $2^{a-k}$ from each of the first $n$ factors. This gives
\begin{align*}
 -\log\Psi(q)
 &=\log2\left(na-\frac{n(n+1)}2\right)
   +\sum_{j=0}^{n-1}\log(1+2^{-j-\theta})
   +\sum_{\ell=1}^\infty\log(1+2^{\theta-\ell})\\
 &=\frac{\log2}{2}(a^2-a+\theta-\theta^2)
   +\sum_{j=0}^{\infty}\log(1+2^{-j-\theta})
   +\sum_{\ell=1}^\infty\log(1+2^{\theta-\ell})
   -\mathcal R(q).
\end{align*}
The removed completion is
$\sum_{j=n}^\infty\log(1+2^{-j-\theta})=\mathcal R(q)$;
this proves \eqref{eq:product-periodic}, also when $n=0$.
For $\theta\in[0,1]$ the two series are uniformly dominated by
geometric series, so their sums are continuous. At zero and one,
their combined values both equal
$\log2+2\sum_{j\geq1}\log(1+2^{-j})$, and the polynomial term
vanishes. Hence the endpoint values agree, giving the periodic
extension. Its boundedness follows from continuity on $[0,1]$.
The termwise derivatives of the two series are bounded by
$(\log2)2^{-j}$ and $(\log2)2^{1-\ell}$, whose sums are $2\log2$
each, and the polynomial term has derivative bounded by
$\tfrac12\log2$; so $\mathcal P$ is continuously differentiable on
$[0,1]$ with $|\mathcal P'|\leq\tfrac92\log2$. A continuous
one-periodic function that is Lipschitz on $[0,1]$ is Lipschitz on
$\R$ with the same constant.

The inequalities $0\leq\log(1+x)\leq x$ give
\eqref{eq:product-remainder}. For $x\geq0$,
$x-x^2/2\leq\log(1+x)\leq x-x^2/2+x^3/3$: both differences vanish
at zero and have derivatives $x^2/(1+x)$ and $x^3/(1+x)$.
Sum these with $x=q^{-1}2^{-r}$; the geometric sums are
$2/q$, $4/(3q^2)$, and $8/(7q^3)$, proving
\eqref{eq:product-refined-remainder}. Neither remainder bound uses
$q\geq1$; that hypothesis enters only through $n\geq0$ in the
decomposition \eqref{eq:product-periodic}.
\end{proof}

In particular the leading term is a negative square of $\log q$,
with a bounded periodic correction. This agrees with the
faster-than-every-power estimate \eqref{eq:product-superpower},
which was sufficient for the scalar dissipation counterexample.

\begin{proposition}[Finite-product certificate]\label{prop:product-certificate}
For $q\geq0$ and integer $k\geq1$, put
\[
 P_k(q)=\prod_{j=1}^k(1+q/2^j)^{-1}.
\]
Then
\begin{equation}
 P_k(q)e^{-q/2^k}\leq\Psi(q)\leq P_k(q).
 \label{eq:product-certificate}
\end{equation}
If $q/2^k<1$, the lower bound
$P_k(q)(1-q/2^k)$ is positive and uses only rational operations
when $q$ is rational. In particular
\begin{equation}
 0.580577558204892402290043892297
       <c_*<0.580577558204892402290043892298.
 \label{eq:overlap-certificate}
\end{equation}
\end{proposition}

\begin{proof}
The omitted log sum lies between zero and
$\sum_{j>k}q2^{-j}=q2^{-k}$. Exponentiating proves
\eqref{eq:product-certificate}; $e^{-x}\geq1-x$ gives its
rational version. At $q=1$, $k=128$, form the exact rational
number $P=\prod_{j=1}^{128}2^j/(2^j+1)$. If $\ell$ and $u$
are the two terminating decimals in \eqref{eq:overlap-certificate},
respectively, integer cross-multiplication verifies
\[
 \ell<1-P\leq c_*\leq1-P(1-2^{-128})<u.
\]
The following complete Python calculation reproduces the finite
arithmetic part with the standard library, without floating point:
\begin{verbatim}
from fractions import Fraction as F
P = F(1)
for k in range(1, 129):
    P *= F(2**k, 2**k + 1)
lo = F("0.580577558204892402290043892297")
hi = F("0.580577558204892402290043892298")
assert lo < 1 - P
assert 1 - P * (1 - F(1, 2**128)) < hi
\end{verbatim}
\end{proof}

The accompanying standard-library script
\texttt{supplement/last\_coagulation\_exact.py} additionally writes
outward-rounded decimal bounds and floating-point evaluations of
\eqref{eq:two-point-counterexample}. Its saved JSON distinguishes
the rational certificate from those illustrative evaluations.
The proofs of asymptotics and scalar obstruction do not depend on
the floating-point examples.
